\begin{theorem}[Tamely ramified odd degree]\label{mod:cases-tameodd}
Let $K/F$ be a cyclic extension of odd prime degree
$\ell=[K:F]$.  Assume that its lower break is exactly $t=0$, that
$f(K/F)=1$, and choose an integral uniformizer $\varpi_K$ which generates
$\mathcal O_K$ over $\mathcal O_F$.  Put
\[
 \varpi_F=N_{K/F}(\varpi_K),\qquad
 \chi_K=\chi_F\circ N_{K/F},\qquad
 \psi_K=\psi_F\circ\operatorname{Tr}_{K/F}.
\]
Thus $\varpi_F$ is a uniformizer.  In each of the four proof-bearing
conductor branches below, Lean proves the literal computational identity
\[
 \Delta_K^{\mathrm{fin}}(\chi_K,\psi_K;\Gamma_K)
 \prod_{\mu\in S(K/F)}
   \Delta_F^{\mathrm{fin}}(\mu,\psi_F;\Gamma_\mu)
 =
 \prod_{\mu\in S(K/F)}
   \Delta_F^{\mathrm{fin}}(\mu\chi_F,\psi_F;\Gamma_{\mu\chi}).
 \tag{*}
\]
Every displayed product is over the literal full finite norm-character
group, including its identity element.  Each denominator in $(*)$ is the
one carried by that actual \texttt{deltaFinite} term; no denominator factor
is suppressed.

If $m_F(\chi_F)=0$ and $n=n_F(\psi_F)$, the tame different and trace
pullback give the exact integer
\[
 n_K(\psi_K)=\ell n+(\ell-1).
\]
The upper conductor-zero formula, with the uniformizer-power admissible
denominator, is
\[
 \Delta_K^{\mathrm{fin}}(\chi_K,\psi_K)
   =\chi_F(\varpi_F)^{\ell n+\ell-1}.
\]
For the identity norm character the lower exponent is $n$, whereas every
nonidentity $\mu$ has exact conductor one, satisfies $\mu(\varpi_F)=1$,
and obeys
\[
 \Delta_F^{\mathrm{fin}}(\mu\chi_F,\psi_F)
 =\chi_F(\varpi_F)^{n+1}
   \Delta_F^{\mathrm{fin}}(\mu,\psi_F).
\]
The implementation retains the identity factor and sums the exact
conductors over the full group:
\[
 n+(\ell-1)(n+1)=\ell n+\ell-1.
\]
This proves $(*)$ without introducing a stationary numerator at conductor
zero.

Suppose next that $m_F(\chi_F)=1$ and that $\chi_F$ is supplied together
with a proof that it is a minimal-conductor member of its $S(K/F)$-orbit.
The public theorem receives actual local character data $\chi_K$ and
$\psi_K$, with the displayed norm and trace pullback equalities; it does not
replace them by freshly chosen data.  Internally Lean constructs, for every
$\mu\in S(K/F)$, the exact norm-character datum $\theta_\mu$ and orbit-twist
datum $\xi_\mu$, satisfying
\[
 \operatorname{char}(\theta_\mu)=\mu,\qquad
 \operatorname{char}(\xi_\mu)=\mu\chi_F.
\]
The tame conductor theorem and orbit minimality then derive, rather than
assume as endpoint-row certificates,
\[
 \begin{aligned}
  &m_K(\chi_K)=1,\qquad m_F(\theta_1)=0,\\
  &m_F(\theta_\mu)=1\quad(\mu\ne1),\qquad
    m_F(\xi_\mu)=1\quad(\text{all }\mu).
 \end{aligned}
\]
In particular a cancelled twist cannot be passed to a conductor-one
residue formula.

The denominators in this branch are also constructed from the actual
characters and conductors.  Set
\[
 \Gamma_F=\varpi_F^{\,1+n_F(\psi_F)},\qquad
 \Gamma_K=\iota_{F,K}(\Gamma_F),
\]
and define
\[
 \Gamma_\mu=\varpi_F^{\,m_F(\theta_\mu)+n_F(\psi_F)},\qquad
 \Gamma_{\mu\chi}
   =\varpi_F^{\,m_F(\xi_\mu)+n_F(\psi_F)}.
\]
Consequently $\Gamma_\mu=\Gamma_F$ for $\mu\ne1$ and
$\Gamma_{\mu\chi}=\Gamma_F$ for every $\mu$, while the identity norm
character keeps its conductor-zero denominator
$\Gamma_1=\varpi_F^{n_F(\psi_F)}$.  The exact trace conductor and the
mapped common-denominator theorem prove that $\Gamma_K$ is admissible for
$(\chi_K,\psi_K)$; in valuation notation this is the equality
\[
 v_K(\Gamma_K)=\ell(1+n_F(\psi_F))
              =1+n_K(\psi_K).
\]
Moreover $\chi_K(\Gamma_K)=\chi_F(\Gamma_F)^\ell$, and for every
nonidentity norm character the common element $\Gamma_F$ is exhibited in
the norm range, so $\mu(\Gamma_F)=1$.

Use the canonical cyclic enumeration
\[
 S(K/F)=\{\mu_j:0\leq j<\ell\},\qquad \mu_0=1,
\]
coming in Lean from \texttt{ramifiedNormCharacterZModEquiv}.  There is no
residual-transport certificate among the public theorem's inputs.  From
$h_{\mathrm{res}}:f(K/F)=1$, Lean proves that the canonical extension
residue map is bijective and packages it as
\[
 e:k_F\simeq k_K.
\]
It then defines
\[
 \tau=\mu_1,\qquad
 \eta=\overline{\theta_\tau},
\]
using the canonical $\operatorname{ZMod}(\ell)$ norm-character
equivalence.  The homomorphism law for that equivalence gives
$\mu_j=\tau^j$; the exact conductor of $\tau$, prime degree, and the
residual-character formula prove internally that
\[
 \operatorname{ord}(\eta)=\ell,\qquad
 \overline{\theta_{\mu_j}}=\eta^j
 \quad(0\leq j<\ell).
\]
In particular the required nonidentity rows for
$j\in\operatorname{Icc}(1,\ell-1)$ and the divisibility
$\ell\mid(|k_F|-1)$ are consequences, not hypotheses.

Likewise the public equalities
$h_{\mathrm{comp}}:\chi_K=\chi_F\circ N$ and
$h_{\mathrm{trace}}:\psi_K=\psi_F\circ\operatorname{Tr}$, together with
the actual mapped denominator pair, are used to prove the upper transports
\[
 \bar\chi_K(e(x))=\bar\chi(x)^\ell,\qquad
 \bar\psi_K(e(x))=\bar\psi(\ell x).
\]
The first follows by lifting a residue unit, applying
$N_{K/F}(\iota(u))=u^\ell$, and reducing again.  The second lifts an
integral representative and applies
$\operatorname{Tr}_{K/F}(\iota(a))=\ell a$ with
$\Gamma_K=\iota(\Gamma_F)$.  Thus the residue equivalence, the canonical
$\eta$ and its order, every norm row, and both upper transports are all
constructed or proved internally.  Here $\bar\chi$ and $\bar\psi$ are
formed from the actual downstairs data and $\Gamma_F$.  Nonvanishing of
the residual additive character, all conductor assertions, equality of
the common denominators, and every local $\Delta$ row below are also
derived; the theorem has no residual certificate parameter.

Applying the conductor-zero and conductor-one formulas from
\texttt{Delta.ResidueFormula} gives, with
$s=\chi_F(\Gamma_F)$,
\begin{align*}
 \Delta_K^{\mathrm{fin}}(\chi_K,\psi_K;\Gamma_K)
  &=-s^\ell\,
    \Aphase\!\left(
      \bar\chi(\ell^\ell)
      \tau_{k_F}(\bar\chi^{\,\ell},\bar\psi)\right),\\
 \Delta_F^{\mathrm{fin}}(\theta_{\mu_j},\psi_F;\Gamma_{\mu_j})
  &=-\Aphase\!\left(\tau_{k_F}(\eta^j,\bar\psi)\right)
       &&\left(j\in\operatorname{Icc}(1,\ell-1)\right),\\
 \Delta_F^{\mathrm{fin}}(\xi_1,\psi_F;\Gamma_{1\chi})
  &=-s\,\Aphase\!\left(\tau_{k_F}(\bar\chi,\bar\psi)\right),\\
 \Delta_F^{\mathrm{fin}}(\xi_{\mu_j},\psi_F;\Gamma_{\mu_j\chi})
  &=-s\,\Aphase\!\left(
      \tau_{k_F}(\bar\chi\eta^j,\bar\psi)\right)
       &&\left(j\in\operatorname{Icc}(1,\ell-1)\right),
\end{align*}
while
$\Delta_F^{\mathrm{fin}}(\theta_1,\psi_F;\Gamma_1)=1$.
The positive factor $\bar\chi(\ell^\ell)$ comes from transporting the
actual upper Gauss sum across $e$ and scaling
$\bar\psi(x)\mapsto\bar\psi(\ell x)$; it is neither inverted nor replaced
by a sign.  The definition of \texttt{langlandsGaussSum} retains the
project's inverse multiplicative-character convention, and every
conductor-one row retains its leading minus sign.

Lean reindexes both full norm-character products through the canonical
$\operatorname{ZMod}(\ell)$ equivalence, separates the identity index, and
applies phases to the exact Hasse--Davenport identity
\[
 \bar\chi(\ell^\ell)
 \tau_{k_F}(\bar\chi^{\,\ell},\bar\psi)
 \prod_{j\in\operatorname{Icc}(1,\ell-1)}
   \tau_{k_F}(\eta^j,\bar\psi)
 =
 \tau_{k_F}(\bar\chi,\bar\psi)
 \prod_{j\in\operatorname{Icc}(1,\ell-1)}
   \tau_{k_F}(\bar\chi\eta^j,\bar\psi).
\]
There is no additional Hasse--Davenport sign: each side of $(*)$ has
exactly $\ell$ conductor-one leading minus signs and the same factor
$s^\ell$.

Finally write a high conductor as
\[
 m_F(\chi_F)=2d+\varepsilon>1,
 \qquad d\geq1,\qquad \varepsilon\in\{0,1\}.
\]
Every norm character has conductor at most one.  Let $\Gamma_F$ be the
explicitly supplied admissible denominator and let
\[
 R:\quad\text{a supplied representative of the actual stationary
 coefficient quotient class for }(\chi_F,\psi_F;\Gamma_F).
\]
Lean transports $R$ to the corresponding Lamprecht lattice, retaining its
proof of quotient-class membership, and constructs the field unit
\[
 \beta=
 \texttt{stableStationaryRepresentativeUnit}(R.\texttt{toLamprecht})
\]
from that supplied representative.  Thus $\beta$ is not an independent or
canonical input.  Exact stable twisting then has the orientation
\[
 \Delta_F^{\mathrm{fin}}(\mu\chi_F,\psi_F;\Gamma_F)
 =\mu(\Gamma_F/\beta)\,
  \Delta_F^{\mathrm{fin}}(\chi_F,\psi_F;\Gamma_F).
\]
The inverse is on $\beta$, and there is no extra sign.  Inversion pairing
in the odd-order norm-character group proves
\[
 \prod_{\mu\in S(K/F)}\mu(\Gamma_F/\beta)=1,
 \qquad
 \prod_{\mu\in S(K/F)}
   \Delta_F^{\mathrm{fin}}(\mu,\psi_F;\Gamma_\mu)=1,
\]
with the identity character retained and with each $\Gamma_\mu$ supplied
as an admissible denominator for its actual norm-character datum.  Hence
the lower twisted product is exactly
\[
 \Delta_F^{\mathrm{fin}}(\chi_F,\psi_F;\Gamma_F)^\ell.
\]

The high comparison uses an actual
\texttt{HighStableOddParameterData} object.  Its ordered denominator pair
has $\Gamma_K$ equal to the image of $\Gamma_F$; the same supplied $R$ is
mapped upstairs.  Norm and trace give the exact $\ell$th power of the
admissible and elementary Lamprecht factors.  If $\varepsilon=0$, supplied
lower and upper even critical coordinates give critical factors equal to
one.  If $\varepsilon=1$, supplied lower and upper odd critical coordinates
and a \texttt{RamifiedHasseComparisonData} certificate in the tame row
$t=0$ prove
\[
 \Delta_{3,K}=\Delta_{3,F}^{\ell}.
\]
This includes the characteristic-two tame normal form when that is the
residue characteristic; odd extension degree does not impose odd residue
characteristic.  In both parity branches one obtains
\[
 \Delta_K^{\mathrm{fin}}(\chi_K,\psi_K;\Gamma_K)
 =\Delta_F^{\mathrm{fin}}(\chi_F,\psi_F;\Gamma_F)^\ell,
\]
and hence $(*)$.  No stationary quotient is represented by a hidden
choice: $R$, its membership proof, the ordered denominator data, the
critical coordinates, and, in the odd branch, the finite Hasse-comparison
certificate remain explicit inputs; $\beta$ is deterministically wrapped
from $R$.

\LeanFile{LanglandsFirstMainLemma/Cases/TameOdd.lean}
\lean{LanglandsFirstMainLemma.firstMain_tameOdd_zero}
\lean{LanglandsFirstMainLemma.firstMain_tameOdd_one}
\lean{LanglandsFirstMainLemma.firstMain_tameOdd_high_even}
\lean{LanglandsFirstMainLemma.firstMain_tameOdd_high_odd}
\lean{LanglandsFirstMainLemma.firstMain_tameOdd}
\leanok
\SourceText{Propositions prop:tame-conductor-one-comparison and prop:tame-odd-new.}
\uses{mod:parameters-ramifiedhassecomparison,mod:finitefield-hassedavenportproduct,mod:lamprecht-stabletwist,mod:delta-residueformula}
\FormalizationContract The principal Lean declaration is a proposition-valued
branch bundle whose fields are the conductor-zero, conductor-one, high-even,
and high-odd computational theorems described above.  It performs neither
minimal-orbit selection nor case dispatch.  A caller selects the appropriate
field and supplies that branch's hypotheses and data: in particular
minimality and the actual norm/trace pullback equalities in the conductor-one
branch, and the actual parameter, representative, critical-coordinate,
denominator, and Hasse data required by the high branch.  The conductor-one
implementation itself constructs the residue-field equivalence, the
canonical residue character $\eta$ and its exact order, every norm-character
residue row, both upper transports, the actual norm and orbit-twist character
data, and all their denominators.  It then derives every needed conductor
equality and every \texttt{deltaFinite} residue row and proves the literal
full-product identity.  In particular there is no residual-transport
certificate argument.  The high implementation constructs $\beta$ from the
explicitly supplied $R$ rather than receiving or choosing a second
representative.  No theorem from \texttt{Cases.Dispatch}, \texttt{Main}, or
another completed case is used.  Passing from these computational branch
theorems to an arbitrary \texttt{LocalConstantFunction} requires the
separate finite--integral bridge and an external minimal-orbit/case
dispatcher.
\end{theorem}
